В настоящей работе получены строгие нижние оценки константы Чватала--Санкоффа $\gamma_\sigma$ для размеров алфавита $\sigma \in \{2, 3, \ldots, 10\}$, превосходящие лучшие опубликованные результаты.

\subsection*{Основные результаты}

\begin{enumerate}
    \item \textbf{Рекорд по бинарному алфавиту:} доказано $\gamma_2 \geq 0{,}7964$. Это улучшение лучшей опубликованной нижней оценки $0{,}79266$~\cite{heineman2024} на $+0{,}0037$ (сопоставимо с суммарным прогрессом по задаче за два десятилетия); при этом сохраняется отрыв $\approx 0{,}030$ от лучшей известной верхней оценки $\gamma_2 \leq 0{,}8263$~\cite{lueker2009}.

    \item \textbf{Рекорды для $\sigma = 3, \ldots, 10$:} получены новые нижние оценки для всех остальных рассмотренных алфавитов; прирост к лучшим известным результатам составляет от $+0{,}17\,\%$ при $\sigma = 3$ до $+5{,}42\,\%$ при $\sigma = 8$.

    \item \textbf{Метод:} разработан и формально обоснован метод оконного динамического программирования. Метод универсально применим для произвольного $\sigma \geq 2$ при помощи единого параметризованного описания состояния и рекурсии.

    \item \textbf{Корректность:} доказаны теоремы~3.3 (корректность ДП-уравнений) и~3.4 (сходимость к величине, не превосходящей $\gamma_\sigma$). Это гарантирует, что все полученные численные значения являются строгими нижними оценками, а не эмпирическими приближениями.

    \item \textbf{Верификация:} все оценки независимо проверены методом Monte-Carlo (классический ДП-LCS на парах случайных строк длины до $20\,000$). Ни одна нижняя оценка не превосходит соответствующего эмпирического среднего, что является необходимым условием корректности.
\end{enumerate}

\subsection*{Авторский вклад}

\begin{itemize}
    \item Сформулирована и обоснована схема стратегии с ограниченным обзором как источника строгих нижних оценок $\gamma_\sigma$.
    \item Разработан и доказан метод динамического программирования вычисления оптимальной такой стратегии (теоремы~3.3, 3.4).
    \item Реализован, оптимизирован и универсализирован алгоритм на C++ (семь итераций оптимизаций x0--x6, специализированная и универсальная версии).
    \item Получены численные рекорды для $\sigma = 2, \ldots, 10$.
    \item Реализована независимая Monte-Carlo верификация.
\end{itemize}

\subsection*{Перспективы дальнейших исследований}

\begin{enumerate}
    \item \emph{Увеличение параметров.} На серверах с большей памятью возможен переход к $N = 13$ и/или $\mathrm{MAX\_SH} = 40$, что даст дальнейшее улучшение оценки. Шаг $N \to N+1$ удваивает потребление памяти и в $\sigma^2 = 4$ раза увеличивает время счёта при $\sigma = 2$, поэтому переход к $N = 13$ потребует $\approx 8$~ГБ и $\approx 20$~часов на одно ядро.
    \item \emph{Параллельные реализации.} Текущая реализация однопоточная; внутренний цикл по парам масок $(a, b)$ тривиально распараллеливается (независимые ячейки слоя $f_K[\mathrm{sh}]$), а внешний цикл по $\mathrm{sh}$ имеет линейную зависимость по предыдущему слою. Грубая оценка: на 16-ядерном сервере ожидается $\approx 12$--$14\times$ ускорение, что позволит ту же оценку получить за десятки минут вместо часов.
    \item \emph{Обобщение на $d > 2$ строк.} Метод естественно обобщается на одновременный поиск LCS $d$ строк; пространство состояний растёт как $O(\mathrm{MAX\_SH}^{d-1} \cdot \sigma^{dN})$. Уже при $d = 3$, $\sigma = 2$ практический предел $N \leq 6$, $\mathrm{MAX\_SH} \leq 15$; интересны константы $\gamma_\sigma^{(d)}$ многомерных задач для $d \in \{3, 4\}$.
    \item \emph{Сокращение зазора.} Разрыв между нашей нижней оценкой $0{,}7964$ и лучшей верхней оценкой $\gamma_2 \leq 0{,}8263$~\cite{lueker2009} составляет $\approx 0{,}030$. Происхождение зазора двояко: с нашей стороны --- ограниченность горизонта обзора стратегии (стратегия с $N = \infty$ совпадает с оптимальной по доказательству); со стороны Lueker --- консервативные мажоранты в потенциальной функции. Дальнейший прогресс возможен как за счёт усиления нижних оценок (увеличение $N$, более выразительные классы стратегий), так и за счёт улучшения верхних.
    \item \emph{Несимметричные модели.} Метод можно адаптировать к моделям с неравными вероятностями символов и к взвешенному LCS, где разные пары символов имеют разные веса совпадения.
    \item \emph{Прикладные следствия.} Точное знание $\gamma_\sigma$ ценно в биоинформатике (оценка значимости выравниваний), в системах сравнения версий (ожидаемая длина общей части двух случайно-возмущённых файлов) и в теории сжатия.
    \item \emph{Открытый теоретический вопрос: монотонность по параметрам.} Эмпирически во всех проведённых расчётах величина $\gamma_{\sigma, N, \mathrm{MAX\_SH}}$ монотонно неубывает по $N$ и $\mathrm{MAX\_SH}$ (см.\,\S5.2 и приложение~Г); формальное доказательство этого факта требует аккуратного согласования информационных состояний стратегий и в данной работе не приводится. Подчеркнём, что данное утверждение \emph{не используется} в обосновании ни одной из полученных численных нижних оценок: каждое отдельное значение является строгой нижней оценкой $\gamma_\sigma$ согласно теореме~3.4 независимо от прочих значений параметров.
\end{enumerate}

Полученные результаты подтверждают эффективность подхода, основанного на анализе локального окна с явным ограничением рассогласования указателей: при умеренных вычислительных ресурсах удалось получить прорыв в задаче, развивавшейся медленно на протяжении полувека.
